%---个人简历、在学期间发表的学术论文及研究成果---
% !TeX root = ../brightPhD.tex

\chapter*{个人简历、在学期间发表的学术论文及研究成果}
\addcontentsline{toc}{chapter}{个人简历、在学期间发表的学术论文及研究成果}
\vspace{1.0cm}
\noindent \heiti{个人简历}

\vspace{0.5em}

\songti{
	\begin{tabular}{p{0.42\textwidth} p{0.42\textwidth}}
		姓名：何亮 & 性别：男 \\
		籍贯：湖南岳阳 & 出生年月：1998年5月 \\
	\end{tabular}
	
	\begin{itemize}
		\item 2016.09--2020.06，在湖南工程学院理学院信息与计算科学专业学习，获理学学士学位。
		\item 2020.09--2023.06，在上海师范大学数理学院运筹学与控制论专业学习，获理学硕士学位。
		\item 2023.09--至今，在湘潭大学数学与计算科学学院数学专业攻读博士学位。
	\end{itemize}
}

\vspace{1pt}
\noindent \heiti{在学期间发表的学术论文}

\vspace{0.5em}

\songti{
	\begin{enumerate}
		\renewcommand{\labelenumi}{[\theenumi]}
		
		\item \textbf{Liang He}, Huayi Wei*, Tian Tian.
		SOPTX: A Modular and Extensible Framework for Topology Optimization with Multi-Backend Support[J].
		Communications in Computational Physics, 已录用。
		
		\item Tian Tian, Chunyu Chen, \textbf{Liang He}, Huayi Wei*.
		Adaptive finite element method for phase field fracture models based on recovery error estimates[J].
		Journal of Computational and Applied Mathematics, 2026, 472: 116732.
		
		\item Yangyang Zheng, Huayi Wei*, Yunqing Huang, Chunyu Chen, Tian Tian, Hanbin Liu, Wenbin Wang, \textbf{Liang He}.
		FEALPy: A Cross-platform Intelligent Numerical Simulation Engine[J].
		Communications in Computational Physics, 已录用。
	\end{enumerate}
}

\vspace{1pt}

\noindent \heiti{其它研究成果}

\vspace{0.5em}

\songti{
	\begin{enumerate}
		\renewcommand{\labelenumi}{[\theenumi]}
		\item 《结构拓扑优化仿真计算软件》，软件著作权，登记号：2025SR0887352。（第一完成人）
		
		\item 《建筑结构力学仿真计算软件》，软件著作权，登记号：2024SR0961048。（第三完成人）
	\end{enumerate}
}


%\songti{
%	 \begin{enumerate}
%	 	\renewcommand{\labelenumi}{[\theenumi]}
%	 	\item 本人为第一作者. SOPTX: A Modular and Extensible Framework for Topology Optimization with Multi-Backend Support[J]. Communications in Computational Physics (已接受).
%	 	
%	 	\item 本人为第三作者. Adaptive finite element method for phase field fracture models based on recovery error estimates[J]. Journal of Computational and Applied Mathematics（）
%	 	\item 本人为第一作者. A FETD scheme and analysis for photonic crystal waveguides comprising third-order nonlinear and linear materials[J]. Journal of Computing Applied Mathematics, 2023, 424, 115005. (SCI 收录)%, WOS:000909190800001, IF=2.4
%	 	\item 本人为第一作者. Total and scattered field decomposition technique in mixed FETD methods and its applications for electromagnetic cloaks[J]. Applied Mathematics Letters, 2024, 153, 109061. (SCI 收录)
%	 	\item 本人为第二作者. FETD study of the wave propagation in chiral metamaterials[J]. Advances in Applied Mathematics and Mechanics, 2021, 13(1), 191--202. (SCI 收录)% WOS:000581925000010, IF=1.4
%	 	\item 本人为第三作者. The material parameter design and finite element simulation of the quadrilateral thermal cloak device[J]. Applied Mathematics Letters, 2019, 94, 99--104. (SCI 收录)%WOS:000465061000015, IF=3.7
%	 \end{enumerate}
%
%}
\vspace{1pt}
%\noindent \heiti{在学期间科研创新项目}
%\songti{	
%\begin{itemize}
%	\item 湖南省研究生科研创新项目(CX2018B380): 求解Dirac方程的高效高精度数值方法.
%\end{itemize}}
